\begin{abstract}
Low-cost, AD5940-based potentiostats are now a recognized instrument class,
with several independently published designs and, as of 2026, a dedicated
review of their architectures. That review treats potentiostat control
software only functionally -- what it does -- and nowhere assesses how it is
built, leaving firmware software-engineering quality unaddressed as a
reliability variable in this literature. We audit the firmware of HunStat2,
our own AD5941/RP2040 potentiostat built directly from the AD5940 datasheet,
standard three-electrode potentiostat design practice, and the Analog
Devices AD5940 evaluation-kit reference firmware, and find it partially
object-oriented in a way that carries real risk: the AFE configuration,
sampling, and current-conversion routine for three of its five electrochemical
methods (chronoamperometry, square-wave voltammetry, differential pulse
voltammetry) is duplicated verbatim three times, a live status-LED bug
silently discards every caller's color argument, and two entire
subsystems -- a command parser and a second electrochemical-methods
implementation, together several hundred lines -- are unreachable, parallel
duplicates of the code paths actually in use. We refactor: the triplicated
routine is extracted into one shared base class inherited by all three
methods; both dead subsystems are removed after tracing every call site to
confirm they are never invoked; the LED bug is fixed. The refactor is not
just argued but measured against real hardware: cyclic voltammetry,
chronoamperometry, square-wave voltammetry, and differential pulse
voltammetry were run against a traceable dummy cell on the physical
instrument before and after the change. Cyclic voltammetry -- the least
noise-sensitive of the four -- shows a mean relative difference of
$0.070\,\%$ and a maximum of $2.83\,\%$ between the two firmware versions
across a $305$-point sweep, consistent with ordinary measurement noise
rather than a behavioral regression; chronoamperometry, square-wave, and
differential pulse voltammetry show the same pattern. A second, independent
real load (a plain $\SI{10}{\kilo\ohm}$ resistor in place of the Randles
dummy cell) reproduces the expected current-voltage relationship on the
refactored firmware alone, corroborating that the change generalizes beyond
one specific load. The refactor also reduced compiled flash usage by
$960$~bytes and removed dead code that a future contributor could otherwise
have mistaken for the live implementation. We conclude that, for this class
of instrument, software architecture quality can be improved substantially
using ordinary refactoring technique -- extract superclass, dead-code
elimination, bug fix -- with a measured, near-zero risk to the measurement
correctness that is the instrument's actual purpose, making it a
no-tradeoff reliability lever that the existing open-source-potentiostat
literature has not evaluated.
\end{abstract}

\begin{keyword}
potentiostat \sep low-cost instrumentation \sep dummy cell \sep
software architecture \sep firmware refactoring \sep AD5940 \sep
open-source hardware
\end{keyword}
